% USER-MAINTAINED SCRATCHPAD.
% Agents must not update this project unless the user explicitly requests an edit to it.
% This is a provisional narrative outline, not a statement of established contributions.
\documentclass[10pt,a4paper]{article}
\usepackage[margin=23mm]{geometry}
\usepackage[T1]{fontenc}
\usepackage{lmodern}
\usepackage{amsmath,amssymb}
\usepackage{titlesec}
\usepackage[hidelinks]{hyperref}
\setlength{\parindent}{0pt}
\setlength{\parskip}{3pt}
\titlespacing*{\section}{0pt}{10pt}{3pt}
\titlespacing*{\subsection}{0pt}{6pt}{2pt}
\titleformat{\section}{\large\bfseries}{\thesection}{.6em}{}
\titleformat{\subsection}{\normalsize\bfseries}{\thesubsection}{.6em}{}
\begin{document}

{\LARGE\bfseries Paper outline: author scratchpad}\par
{\small 25 October 2026 \quad|\quad Provisional narrative, before results}\par
{\footnotesize Personal notes: agents should edit only on explicit request.
Literature review and related work are assumed; proposed claims and decoder versions remain open.}

\section{System model and computational problem}
Define the unsourced Gaussian multiple-access model, known activity/noise assumptions, message multiplicities and multiset PUPE, then explain how $B$, $M=2^B$, $K$ and blocklength determine the computational challenge.

\section{Structured polar transmission design}

\subsection{Why polar coding and transmission patterns?}
Motivate the algebraic structure and compact encoder description, with optional dense--polar comparisons at $B=12/14$ presented as receiver-dependent evidence rather than proof of equal encoding performance.

\subsection{Deterministic energy balancing}
Derive resource-energy balancing from the second-moment geometry, distinguishing its interference criterion from PUPE optimality and leaving room for a better encoder construction.

\subsection{What structure does the decoder gain?}
Identify useful pattern, amplitude and count relations; keep the intuition of one or two extra bits as something to quantify, rather than assuming additional payload or independent information.

\section{The exact decoding target}

\subsection{Posterior and decision objective}
Express the physical likelihood and source prior through the polar construction, then derive the count-tail decision that minimizes expected multiset PUPE.

\subsection{Useful exact reductions}
Select the exact observations needed by the eventual algorithm---such as permutation symmetry and polar factorization---and distinguish representation savings from the unchanged $M$-message alphabet and combinatorial count-state search.

\section{From the exact target to a decoder}
Present only approximations that contribute to the final method, explaining the information each discards, its error control and its computational cost; the following version sequence is provisional.

\subsection{V1: core nonlearned decoder}
Build the basic inference method around the strongest discovery mechanism, potentially multiset beam search or guided importance sampling.

\subsection{V2: complementary discovery or correction}
Add a justified way to recover hypotheses or posterior mass missed by the core method, with correct weighting and count readout.

\subsection{V3: targeted refinement}
Introduce useful local or joint corrections, showing through controlled comparisons what they add beyond V1/V2 and whether their cost is justified.

\subsection{Optional V4: learned internal approximations}
If pursued later, learn to accelerate or improve specific internal computations while retaining the algebraic target and a clear nonlearned reference.

\section{Evaluation methodology}

\subsection{A tractable finite-alphabet test regime}
Justify $B=12/14$ for explicit-codebook comparisons and test the proposed low-collision analogy to larger $B$, separating message collisions from pattern overlap and reserving exhaustive optimal decoding for genuinely tiny cases.

\subsection{Matched baselines, versions and ablations}
Compare PUPE and computational cost on common inputs against D2, the appropriately labelled ODMA baseline and relevant bounds, using tiny exact references and ablations to isolate each contribution.

\section{Results}
\emph{Reserved.}

\end{document}
